\documentclass[aspectratio=169]{beamer}
\usepackage{ctex, hyperref}
\usepackage[T1]{fontenc}

% other packages
\usepackage{latexsym,amsmath,xcolor,multicol,booktabs}
\usepackage{graphicx,pstricks,listings,stackengine}

% ===== 可编辑：报告信息 =====
\author{李沁轩}
\title{CUMT Beamer Theme}
\subtitle{毕业设计开题报告}
\institute[计算机科学与技术学院/人工智能学院]{计算机科学与技术学院/人工智能学院\\计算机科学与技术24-3班}
\date{2026年9月27日}
% ===== 报告信息结束 =====
\usepackage{CUMT}
\usepackage{CUMT-campus}

% 可编辑：致谢文案，版式由主题生成。
\renewcommand{\CUMTThanksEnglish}{Thanks!}
\renewcommand{\CUMTThanksChinese}{恳请各位评委批评指正}

% defs
\def\cmd#1{\texttt{\color{red}\footnotesize $\backslash$#1}}
\def\env#1{\texttt{\color{blue}\footnotesize #1}}
\definecolor{deepblue}{rgb}{0,0,0.5}
\definecolor{deepred}{rgb}{0.6,0,0}
\definecolor{deepgreen}{rgb}{0,0.5,0}
\definecolor{halfgray}{gray}{0.55}

\lstset{
    basicstyle=\ttfamily\small,
    keywordstyle=\bfseries\color{CUMT},
    emphstyle=\ttfamily\color{deepred},    % Custom highlighting style
    stringstyle=\color{deepgreen},
    numbers=left,
    numberstyle=\small\color{halfgray},
    rulecolor=\color{CUMTLine},
    backgroundcolor=\color{CUMTPale},
    frame=single,
}


\begin{document}

\heiti
% 封面由主题生成；只需修改导言区的标题、作者等信息。
\CUMTMakeTitle

\begin{frame}{目录}
    \tableofcontents[sectionstyle=show,subsectionstyle=show/shaded/hide,subsubsectionstyle=show/shaded/hide]
\end{frame}


% 可编辑内容层：替换章节、标题和正文；保留普通 frame 环境。
% 章节决定顶部导航的名称，页面决定导航圆点。
\section{项目背景}
\begin{frame}{研究背景}
  \begin{itemize}
    \item 在此介绍问题背景与研究意义。
    \item 在此说明现有方法与待解决的问题。
  \end{itemize}
\end{frame}

\section{研究内容}
\begin{frame}{研究目标}
  \begin{block}{核心目标}
    在此填写研究目标、技术路线或主要贡献。
  \end{block}
\end{frame}

\section{工作计划}
\begin{frame}{后续安排}
  \begin{enumerate}
    \item 文献调研与方案设计。
    \item 方法实现与实验验证。
    \item 总结成果与撰写论文。
  \end{enumerate}
\end{frame}


% 如需替换致谢文案，在导言区重定义 \CUMTThanksEnglish / \CUMTThanksChinese。
\CUMTMakeThanks

\end{document}
